% !TeX root = ../papers/nc-p3.tex
\section{種数5の四つの分岐値における関係格子}
\label{sec:rank-fifteen-pencil-ramified-ends}
\begin{theorem}\label{thm:rank-fifteen-pencil-ramified-ends}
一般の種数5水準の両表示で、$K=\pm s\pm1$ 上の16境界点のそれぞれに、
通常窓が標準Koszul左右AS正則になる開近傍がある。さらに $c$ が101で
$3/4$ に還元され、$K$ がこの四値のいずれかに還元される全ての通常な標数零窓は
同じ正則性を持ち、各始点で四次中心の多項式環上階数256の自由加群となる。
一般水準の56境界点のうち、これまでの近傍判定で残るのは $K=\pm1$ 上の八点である。
\end{theorem}
\begin{proof}
$\delta,\epsilon\in\{1,-1\}$ とし、
$K_0=-\delta s-\epsilon$ と置く。$x$ は $\delta=1$ で $r$、
$\delta=-1$ で $w$ とし、その半トレースを $E=K+\delta(K^2-1)/s$ とする。
$E(K_0)=\epsilon$、$x=\epsilon$ である。他方の根を $q$ とすると境界値では
\[
 q^2-2(2K_0-\epsilon)q+1=0,\qquad n_0=\pm(c+\delta\epsilon)/c .
\]
前式の二根は相異なり、$E'(K_0)\ne0$、$n_0\ne0$ である。
これらは指定する101上の点でも単元条件を満たす。
$u=x-\epsilon$ と取ると $(x+x^{-1})/2=\epsilon+O(u^2)$ であり、
$c$ を係数に保って $K=K_0+O(u^2)$、$n=n_0+O(u^2)$、
他方の傾きも $q+O(u^2)$ と展開する。従って一次項の計算では $x$ だけを動かせばよい。

各位置・文字の関係行列から全最大小行列式を取り、共通分母と共通因子を除いて
$z,\lambda,a,b,e,f$ の多項式の組を得る。ここに各表示の窓を代入する。
\path{computations/verify_rank_fifteen_pencil_ramified_ends.py} は $c=3/4$ を
101で特殊化し、各 $\delta,\epsilon$ と両表示について $n_0,q$ の全四択を取る。
各組の定数項に単元があればその組を使う。定数項が全て零に還元される場合は、
$\mathbb Q(c)[q]$ で上記のモニック二次式を法として各分子の剰余が
\emph{恒等的に零}になること、分母が同二次式と互いに素であることを確かめる。
そのため $u$ で一度割ることが可能である。用いた窓の式の分母は $c,r,w$ の積の約数で、
これらは剰余でも単元である。割った後の定数項は $x$ による一次微分であり、
実際に各組で少なくとも一座標が単元となる。これにより、$c$ も変動する近傍で
元の関係を表すGrassmann座標が得られる。

各組の単元座標から行基底を復元すると、各剰余で二次関係の階数6、
経路次元35,56、中心の文字次元 $(3,2,0,0)$、全四始点の長さ13までの
四中心商の標準列を得る。前定理と同じ持上げにより、通常窓の近傍正則性が従う。
境界の有限被覆上で $c=3/4$ の全選択を認証したので、一般水準の全16点に適用できる。
無限遠八点と有限非分岐24点を合わせて48点となる。
\end{proof}

\begin{proposition}
残る $K=k=\pm1$ の八点では、各分岐と表示を固定すると、二次関係の
Grassmann極限は $c$ に依存しない。その二次代数の経路次元は
全始点で長さ3,4,5の順に $20,38,72$ である。
\end{proposition}
\begin{proof}
局所座標に $n$ を取る。$K^2=1+(1+c^2)n^2$、$R=K+cn^2$、$W=K-cn^2$ より
\[
 K=k+O(n^2),\quad r=k+\sigma(c+k)n+O(n^2),\quad
 w=k+\tau(c-k)n+O(n^2),\qquad \sigma,\tau\in\{1,-1\}.
\]
実際 $(r+r^{-1})/2=R$、$(w+w^{-1})/2=W$ の二次項を比べればよい。
両表示で $(z,\lambda,a,b,e,f)$ の境界値は全て1である。
上の検証コードは同じ多項式座標をこの点で評価し、組の定数項が全て零ならば、
表示の一次係数に沿う方向微分を取る。各組は一次までに非零となる。
非零ピボットで割った全係数を $\mathbb Q(c)$ で約分すると、全16組で分子・
分母の $c$ 次数はともに零である。この有理数の関係行列から長さ5まで道を延ばすと
上記の次元を得る。通常窓の四次・五次経路を復元するには、ここでも高次関係の飽和が要る。
\end{proof}
